\documentclass[12pt]{article}
\usepackage[utf8]{inputenc}
\usepackage[T1]{fontenc}
\usepackage{hyperref}
\title{Gudang Language Resources for K-12 Curriculum Development}
\author{Compiled by Claude Sonnet 5.5 (claude-sonnet-5-5) for Andr\'as Kornai}
\date{}
\begin{document}
\maketitle

\section{Introduction}
Gudang (also Djagaraga; ISO 639-3 xgd) is the traditional language of the Gudang people. It is the northernmost language of Cape York Peninsula. Austlang records no known speakers and very little language data. The surviving record is mostly a wordlist collected by Frank Jardine at Somerset, published in Curr's \emph{The Australian Race} (1886/1887), and a comparative Kowrarega--Gudang list in the appendix to the \emph{Narrative of the Voyage of HMS Rattlesnake} (1852).

\section{Primary Educational Resources}
No primer or textbook was found. The Pama Language Centre is reported to be supporting descendants of Gudang people in language revival, but no teaching materials were seen.

\section{Linguistic and Reference Materials}
\begin{itemize}
\item Wikipedia, ``Gudang language'': \url{https://en.wikipedia.org/wiki/Gudang_language}
\item AIATSIS Austlang record Y191: \url{https://collection.aiatsis.gov.au/austlang/language/Y191}
\item Glottolog reference for Jardine's Gudang wordlist (1886): \url{https://glottolog.org/resource/reference/id/151737}
\item Glottolog reference for Harper (1992): \url{https://glottolog.org/resource/reference/id/3239}. Its content was not verified.
\item State Library of Queensland blog posts on the language: \url{https://www.slq.qld.gov.au/blog/2019-international-year-indigenous-languages-word-week-week-eleven} and \url{https://www.slq.qld.gov.au/blog/language-week-week-twenty-djagaraga}. One gives the word \emph{aigi}, ``sandy beach''.
\end{itemize}

\section{Religious and Community Resources}
\begin{itemize}
\item Pama Language Centre, which is reported to support Gudang revival: \url{https://www.pamacentre.org.au/?p=987}
\end{itemize}
No scripture translations were found.

\section{Digital and Audio Resources}
\begin{itemize}
\item International Year of Indigenous Languages page: \url{https://en.iyil2019.org/?p=14011}. I did not verify that it contains Gudang material.
\end{itemize}
No audio or software resources were found.

\section{Curriculum Development Context}
The language record is very small, and the full text of the two historical wordlists was not seen. Any curriculum would have to be built with Gudang descendants and the Pama Language Centre. It could use the few known words, place names and history, and reconstruction would be limited.
\end{document}